% SPDX-FileCopyrightText: © 2026 Idiap Research Institute <contact@idiap.ch>
%
% SPDX-FileContributor: Andre Freitas <andre.freitas@idiap.ch>
% Neuro-symbolic AI Group
%
% SPDX-FileContributor: Danilo Gusicuma <danilo.gusicuma@idiap.ch>
% Neuro-symbolic AI Group
%
% SPDX-License-Identifier: MIT
\subsection{Target navigation: choose an obligation before doing research}
\label{app:granular-target-navigation-loop}
The linked dependency-cartographer skill supplies the conceptual procedure hidden
by this compact boundary. It reads the target's acceptance conditions against a
recursive dependency graph, distinguishes a missing premise from a bridge object,
route lemma, proof debt or stale route, and groups unresolved nodes into coherent
proof tasks. Priority is judged by how a node serves the target, not merely by
its local score. Similarity does not count as logical dependence, and a mentioned
lemma does not count as a discharged obligation. Cycle detection prevents a route
from supporting itself through a chain of references. The crucial inference is
from the \emph{structure of what is missing} to a justified next research question.
This is a prescribed interpretation procedure linked to the workflow, rather than
an algorithm encoded by the native projection signature.

Target navigation separates identifying what remains to be established from
establishing it. Write $T$ for \texttt{Unscoped\allowbreak Research\allowbreak Target}, $F$ for
\texttt{Target\allowbreak Dependency\allowbreak Frontier}, and $N$ for
\texttt{Target\allowbreak Navigation\allowbreak Disposition}. The native occurrences are an unresolved
projection $P_T$ and an evidential assessment $E_T$:
\[
\begin{aligned}
&P_T:T\to\operatorname{Generated}[F],\qquad
 E_T:F\to\operatorname{Generated}[N],\\
&f=P_T(t),\qquad n=E_T(f.\mathrm{value}),\\
&N=\{\mathrm{WorkReady},\mathrm{ReadyForReleasePromotion},\mathrm{Stalled}\}.
\end{aligned}
\]
The first occurrence, \texttt{dependency\_\allowbreak frontier}, changes an identified target
into a target with a referenced dependency frontier. It makes the unresolved
obligations available for selection, but adds no warrant for the target.
Crucially, the original \texttt{pre:\_\allowbreak } does not encode the linked procedure:
$P_T$ names a declared input/output boundary, not a recovered cartography
algorithm. The second occurrence, \texttt{target\_\allowbreak disposition}, assesses that
frontier to choose the next kind of work. The relation
\texttt{Dependency\allowbreak Frontier\allowbreak Enables\allowbreak Target\allowbreak Disposition} records the declared dependency
of the generated decision on the frontier; its generic source/target certificate
shape does not itself establish that the frontier is complete. The displayed
$\mathrm{value}$ selection is the native explicit wrapper adapter.

For a constructed example, suppose the target is a decay bound uniform over a
family of networks. A frontier can distinguish the proved conditional estimate
from the open obligation that the spectral gaps have a positive common lower
bound. Since that obligation remains open, \texttt{Work\allowbreak Ready} directs work to the
gap question; it cannot mean that uniform decay has been proved. This illustrates
why selecting an obligation is a distinct conceptual move from deriving a result.
The native function returns $n$. The original control policy instead switches
work-ready to \texttt{abducer\_\allowbreak work\_\allowbreak loop}, release-candidate to
\texttt{release\_\allowbreak promotion\_\allowbreak loop}, and stops on stall. Those transfers are not
executed by the native return statement, and the native declarations do not encode the skill-level
judgment that chooses among them.

% SPDX-FileCopyrightText: © 2026 Idiap Research Institute <contact@idiap.ch>
%
% SPDX-FileContributor: Andre Freitas <andre.freitas@idiap.ch>
% Neuro-symbolic AI Group
%
% SPDX-FileContributor: Danilo Gusicuma <danilo.gusicuma@idiap.ch>
% Neuro-symbolic AI Group
%
% SPDX-License-Identifier: MIT
% Generated from step-explanations.json
\paragraph{Step-by-step reconstruction.}
The target is a diffusion bound uniform across a graph family. A conditional decay theorem is available, but its common spectral-gap premise is open.
The following is a source-guided illustrative reconstruction. Internal procedures are prescriptions; the native declarations do not establish their execution.
\begin{enumerate}[leftmargin=*,label=\arabic*.]
\item \textbf{Construct the dependency frontier} (\texttt{dependency\_\allowbreak{}frontier}).
\begin{description}[style=standard,labelindent=0pt,leftmargin=5.5em,labelwidth=5em,labelsep=0.5em,align=left,font=\normalfont\bfseries,itemsep=0.15em,parsep=0pt,topsep=0.3em]
\item[Purpose] Locate the target’s remaining proof obligations.
\item[Input] Target acceptance conditions; recursive dependency graph
\item[Method] The native producer is unresolved. The linked cartographer prescribes classifying missing premises, bridge objects, route lemmas and proof debt, while checking cycles.
\item[Check] A dependency must support the target logically; similarity or a mentioned lemma does not discharge it.
\item[Output] Intended frontier: open uniform-gap proof obligation
\item[Next] Disposition assessment requires this frontier; the signature does not establish that cartography was performed.
\item[Example] Separate the available conditional decay theorem from the unproved assertion that every graph has gap at least the same c greater than zero.
\end{description}
\item \textbf{Select a justified obligation} (\texttt{target\_\allowbreak{}disposition}).
\begin{description}[style=standard,labelindent=0pt,leftmargin=5.5em,labelwidth=5em,labelsep=0.5em,align=left,font=\normalfont\bfseries,itemsep=0.15em,parsep=0pt,topsep=0.3em]
\item[Purpose] Select work that controls target completion.
\item[Input] Dependency frontier, including unresolved support
\item[Method] Assess which open node controls the target and distinguish work readiness from release readiness or stall.
\item[Check] An open necessary premise prevents treating the target as complete.
\item[Output] Work disposition: common-gap task and justification
\item[Next] Original policy routes onward; the native return neither proves the premise nor executes that route.
\item[Example] Because uniform decay depends on the open common-gap premise, select that premise for investigation rather than re-deriving the conditional theorem.
\end{description}
\end{enumerate}
\paragraph{Changed-input contrast.}
If the common-gap premise is discharged for the target family, navigation should reassess the remaining acceptance obligations. It must not continue selecting the same already closed gap merely because its node is prominent.


\subsection{Release promotion: conserve the scope of the evidence}
\label{app:granular-release-promotion-loop}
The release-promoter skill prescribes a comparative interpretation of evidence:
separate robust support from missing or fragile support, inspect whether the
active route still addresses the target, and determine whether the chosen
mathematical machinery has the objects it needs. This last question can change the
research decision even when every report validates: endpoint spectral machinery
cannot replace an absent construction or transfer object. A route that repeatedly
hits the same blocker must explain what was learned and what would trigger a
pivot. Promotion consequently depends on whether the proposed claim survives
these comparisons; repair can change the target scope or conceptual route rather
than merely regenerate a report. These are skill-level prescriptions associated
with the release contract, not extra native occurrences or evidence of execution.

A release decision bounds the result supported by a research iteration.
Here $V$ denotes \texttt{Promotion\allowbreak Evidence\allowbreak Packet}, $Q$ denote
\texttt{Claim\allowbreak Bounded\allowbreak Release\allowbreak Proposal}, and $D$ denote
\texttt{Release\allowbreak Promotion\allowbreak Disposition}. Its typed boundary is
\[
\begin{aligned}
&P_R:V\to\operatorname{Generated}[Q],\qquad
 E_R:Q\to\operatorname{Generated}[D],\\
&q=P_R(v),\qquad d=E_R(q.\mathrm{value}),\\
&D=\{\mathrm{Promoted},\mathrm{RepairTarget},
       \mathrm{RepairConcept},\mathrm{Rejected}\}.
\end{aligned}
\]
At \texttt{release\_\allowbreak proposal}, the evidence packet's run and evidence references
are projected into a proposal with run-identity and proposal-reference fields.
Preserving run identity is intended, but equality of the two string fields is
not enforced by these nominal types. Conceptually this is the move from accumulated evidence to a bounded
claim about it. The source's \texttt{pre:\_\allowbreak } leaves proposal construction open;
\texttt{Unresolved\allowbreak Release\allowbreak Proposal\allowbreak Projection} must not be read as an implemented
claim-selection algorithm. At \texttt{promotion\_\allowbreak disposition}, evidential
assessment asks whether that proposal is supportable, needs a target or conceptual
repair, or should be rejected. \texttt{Release\allowbreak Proposal\allowbreak Enables\allowbreak Promotion\allowbreak Disposition}
licenses the declared proposal-to-assessment dependency. It supplies neither an
independent proof nor a rule that reference-bearing proposals must pass. There is
no native multi-input join: the assessment receives the unwrapped proposal, whose
reference must carry the evidence boundary.

For a constructed example, a Yang--Mills route has a finite-regulator correlation
estimate and proposes promoting it as evidence that its continuum construction is
complete. The release comparison exposes an absent same-theory transfer and
nonuniform error control. Repairing the target can license a finite-result claim;
repairing the concept instead asks for the missing transfer object. Passing the
workflow validators cannot decide between those scientific scopes. The skill's
machinery assessment makes the difference explicit: a spectral conclusion cannot
be obtained by invoking endpoint machinery before its analytic preconditions are
available.
The original policy accepts a promoted outcome, routes target repair to target
navigation and conceptual repair to conceptual hypothesis exploration, and stops
on rejection. The native function returns the wrapped disposition and contains
none of those switches. Thus the workflow's contribution is scope calibration,
not an additional experiment or an automatic truth certificate.

% SPDX-FileCopyrightText: © 2026 Idiap Research Institute <contact@idiap.ch>
%
% SPDX-FileContributor: Andre Freitas <andre.freitas@idiap.ch>
% Neuro-symbolic AI Group
%
% SPDX-FileContributor: Danilo Gusicuma <danilo.gusicuma@idiap.ch>
% Neuro-symbolic AI Group
%
% SPDX-License-Identifier: MIT
% Generated from step-explanations.json
\paragraph{Step-by-step reconstruction.}
A release packet contains a finite-regulator correlation estimate but proposes a continuum conclusion without a same-theory transfer or uniform error bound.
The following is a source-guided illustrative reconstruction. Internal procedures are prescriptions; the native declarations do not establish their execution.
\begin{enumerate}[leftmargin=*,label=\arabic*.]
\item \textbf{Construct a bounded proposal} (\texttt{release\_\allowbreak{}proposal}).
\begin{description}[style=standard,labelindent=0pt,leftmargin=5.5em,labelwidth=5em,labelsep=0.5em,align=left,font=\normalfont\bfseries,itemsep=0.15em,parsep=0pt,topsep=0.3em]
\item[Purpose] Bound a release claim by accumulated support.
\item[Input] Iteration evidence packet; target identity; support references
\item[Method] The native construction is unresolved. The linked promoter prescribes separating robust support from fragile support and exposing target drift and missing machinery prerequisites.
\item[Check] The proposed scope cannot exceed the regime actually supported by evidence.
\item[Output] Intended proposal: supported finite estimate; unsupported continuum extension
\item[Next] Promotion assessment consumes the proposal; nominal run references do not enforce identity preservation.
\item[Example] Write the finite-regulator estimate as supported within its assumptions and list continuum passage and uniform error control as open obligations.
\end{description}
\item \textbf{Assess promotion or repair} (\texttt{promotion\_\allowbreak{}disposition}).
\begin{description}[style=standard,labelindent=0pt,leftmargin=5.5em,labelwidth=5em,labelsep=0.5em,align=left,font=\normalfont\bfseries,itemsep=0.15em,parsep=0pt,topsep=0.3em]
\item[Purpose] Assess promotion against the evidence boundary.
\item[Input] Release proposal; underlying support references
\item[Method] Compare claim scope with missing construction and transfer objects; distinguish target repair, conceptual repair and rejection.
\item[Check] Endpoint spectral machinery cannot replace an absent analytic transfer premise.
\item[Output] Promotion disposition: bounded promotion or specified repair
\item[Next] Original policy may promote, route repair or stop; native output alone performs none of these effects.
\item[Example] Restrict the release to the finite estimate or request the missing transfer object; do not accept the continuum announcement from the same packet.
\end{description}
\end{enumerate}
\paragraph{Changed-input contrast.}
Restricting the proposal to the evidenced finite-regulator result can make a bounded claim supportable. Adding more fixed-regulator examples cannot by itself repair the missing uniform transfer.


\subsection{Timeline-to-paper: translate a reasoning record without upgrading it}
\label{app:granular-timeline-paper-independent-workflow}
The linked argument-architect and failed-route-interpreter skills explain how
the record becomes a scientific argument. The architect chooses the paper's form
from the evidence: theorem, worked calculation, obstruction, method, negative
result or positioning note. It identifies the smallest defensible contribution
and maps it to the evidence, novelty boundary and remaining validation burden.
The failure interpreter compresses a route into its original idea, the pressure
that defeated it, the learned invariant and the narrowed object that survives.
Their conceptual product is an explanation of why the current claim is the right
claim, rather than a chronological list of attempts. Section responsibilities
then distribute that argument across the paper. The skill procedures specify this
reasoning even though the native claim-set projection does not implement it;
manifest membership alone supplies no execution order between the linked skills.

Paper construction has an epistemic boundary: reorganizing a reasoning record
must not turn its open conjectures into established results. Let $L$ be
\texttt{Timeline\allowbreak Reasoning\allowbreak Outputs}, $C$ be
\texttt{Bounded\allowbreak Manuscript\allowbreak Claim\allowbreak Set}, and $D$ be
\texttt{Timeline\allowbreak Paper\allowbreak Disposition}. The two declared occurrences give
\[
\begin{aligned}
&P_P:L\to\operatorname{Generated}[C],\qquad
 E_P:C\to\operatorname{Generated}[D],\\
&c=P_P(\ell),\qquad d=E_P(c.\mathrm{value}),\\
&D=\{\mathrm{PaperReady},\mathrm{ClaimBoundaryFailure},\mathrm{Stalled}\}.
\end{aligned}
\]
The \texttt{bounded\_\allowbreak claim\_\allowbreak set} occurrence translates a timeline reference into
a source reference and claim-ledger reference. This establishes presentation
input, not new evidence. Its unresolved projection corresponds to the original
\texttt{pre:\_\allowbreak }; the declarations do not specify how competing claims are
selected, reconciled, or written. The \texttt{paper\_\allowbreak disposition} occurrence
assesses whether this bounded set supports paper assembly.
\texttt{Claim\allowbreak Set\allowbreak Enables\allowbreak Paper\allowbreak Disposition} records the assessment's dependence on
that claim set, with explicit unwrapping at $c.\mathrm{value}$. It does not assert
that the ledger faithfully reflects its source merely because both references
exist. Neither types nor topology specify an internal drafting sequence.

Consider a constructed diffusion study whose timeline contains a conditional
decay derivation, an unsuccessful attempt to generalize it, and the cycle-graph
counterexample to a size-independent spectral gap. A faithful claim set separates
these three statuses. The paper can explain why the generalization fails while
retaining the conditional theorem. If a draft instead states uniform decay for
all connected graphs, \texttt{Claim\allowbreak Boundary\allowbreak Failure} is an appropriate assessment:
the presentation has dropped precisely the assumption challenged by the record.
The outcome is a corrected scientific claim boundary, not merely a formatted
manuscript. The original workflow accepts paper-ready and stops on claim-boundary
failure or stall. The native function only returns $d$; its readiness value is
not itself a publication act. This distinction matters because readable exposition
can improve understanding without increasing the evidential status of any claim.

% SPDX-FileCopyrightText: © 2026 Idiap Research Institute <contact@idiap.ch>
%
% SPDX-FileContributor: Andre Freitas <andre.freitas@idiap.ch>
% Neuro-symbolic AI Group
%
% SPDX-FileContributor: Danilo Gusicuma <danilo.gusicuma@idiap.ch>
% Neuro-symbolic AI Group
%
% SPDX-License-Identifier: MIT
% Generated from step-explanations.json
\paragraph{Step-by-step reconstruction.}
A diffusion timeline contains a conditional decay derivation, a failed uniform generalization and a cycle-graph counterexample to a common spectral gap.
The following is a source-guided illustrative reconstruction. Internal procedures are prescriptions; the native declarations do not establish their execution.
\begin{enumerate}[leftmargin=*,label=\arabic*.]
\item \textbf{Form the defensible claim set} (\texttt{bounded\_\allowbreak{}claim\_\allowbreak{}set}).
\begin{description}[style=standard,labelindent=0pt,leftmargin=5.5em,labelwidth=5em,labelsep=0.5em,align=left,font=\normalfont\bfseries,itemsep=0.15em,parsep=0pt,topsep=0.3em]
\item[Purpose] Recover a supported argument from the research record.
\item[Input] Timeline claims; premises; counterexamples; unresolved attempts
\item[Method] The native projection is unresolved. The argument architect selects the smallest supported contribution; failure interpretation extracts the defeated idea, learned obstruction and surviving statement.
\item[Check] The paper mode must follow the evidence rather than the ambition of the original target.
\item[Output] Intended claim ledger: conditional theorem; refuted generalization
\item[Next] Paper assessment consumes this bounded ledger; the native boundary supplies no automatic drafting algorithm.
\item[Example] Preserve the conditional spectral estimate, report cycles as the obstruction to uniformity, and explain why the unrestricted claim was abandoned.
\end{description}
\item \textbf{Check presentation readiness} (\texttt{paper\_\allowbreak{}disposition}).
\begin{description}[style=standard,labelindent=0pt,leftmargin=5.5em,labelwidth=5em,labelsep=0.5em,align=left,font=\normalfont\bfseries,itemsep=0.15em,parsep=0pt,topsep=0.3em]
\item[Purpose] Check that manuscript organization preserves claim scope.
\item[Input] Bounded claim set; source references
\item[Method] Compare proposed manuscript claims with the premises and failure record, locating any omitted assumption or contradictory status.
\item[Check] A claim-ledger reference is insufficient unless its content faithfully represents the record.
\item[Output] Paper disposition: readiness or specified claim-boundary failure
\item[Next] Original policy accepts readiness or stops; native return is not publication.
\item[Example] A draft asserting uniform decay for all connected graphs fails because it removes the very common-gap premise challenged by the cycle example.
\end{description}
\end{enumerate}
\paragraph{Changed-input contrast.}
Removing the counterexample from the narrative while retaining the broad claim breaks fidelity to the record. Restoring the failed premise and scoped conditional theorem repairs the claim boundary without creating new evidence.


\subsection{Yang--Mills overlay: keep representation and continuum claims apart}
\label{app:granular-ym-representational-dynamics-overlay}
The linked Yang--Mills skills describe a much richer scientific program behind
this audit boundary. Observable-core construction asks whether the chosen
observables separate states and generate enough physical vectors to see every
relevant spectral sector. Representation construction independently specifies
source and target carriers, observables, states and dynamics, then asks what an
observable pullback and state pushforward preserve. Exact equivalence is
separated from coarse-graining and approximation with explicit error dependence.
Continuum coherence asks whether regulated correlations converge in a topology
strong enough to preserve positivity and spectral estimates, with justified limit
orders and theory identity. Rigidity transfer then asks whether a uniform decay
estimate extends from a dense physical class to spectral exclusion. These are
prescribed mathematical burdens, not solved theorems or an execution sequence
inferred from the skill list. The native two-stage audit compresses their results;
it does not make the underlying conceptual questions unknown.

This overlay tests whether a representation candidate stays within the declared
continuum boundary. Write $Y$ for \texttt{Yang\allowbreak Mills\allowbreak Target\allowbreak Context}, $A$ for
\texttt{Continuum\allowbreak Boundary\allowbreak Audit\allowbreak Input}, and $D$ for
\texttt{Yang\allowbreak Mills\allowbreak Overlay\allowbreak Disposition}:
\[
\begin{aligned}
&P_Y:Y\to\operatorname{Generated}[A],\qquad
 E_Y:A\to\operatorname{Generated}[D],\\
&a=P_Y(y),\qquad d=E_Y(a.\mathrm{value}),\\
&D=\{\mathrm{Accepted},\mathrm{Quarantined},\mathrm{Stalled}\}.
\end{aligned}
\]
At \texttt{continuum\_\allowbreak audit\_\allowbreak input}, the projection supplies target-identity and
continuum-boundary fields together with a representation-candidate reference.
Those fields can retain the input context, but preserving their values is an
obligation not enforced by the nominal string types.
Its conceptual purpose is to put a candidate beside the boundary it must respect.
It supplies no bridge from a finite representation to a continuum assertion.
The original \texttt{pre:\_\allowbreak } leaves that preparation internally unspecified;
one cannot infer literature search, simulation, or theorem construction as extra
native occurrences. At \texttt{overlay\_\allowbreak disposition}, evidential assessment
judges the prepared candidate against the boundary.
\texttt{Audit\allowbreak Input\allowbreak Enables\allowbreak Overlay\allowbreak Disposition} connects this audit input to its
generated assessment, without certifying that the boundary has been crossed or
that an accepted representation proves the scientific target.

For a constructed illustration, suppose a selected observable has exponentially
decaying correlations. Add an orthogonal sector unseen by that observable whose
positive energies approach zero. The selected correlation is unchanged, yet a
global spectral-gap claim is defeated. Observable-core reasoning identifies the
missing density obligation; rigidity transfer restricts the conclusion to the
visible sector until that obligation is met. This countermodel distinguishes
``the measurement decays'' from ``all physical excitations are gapped.'' It is a
test of the prescribed inference, not a proposed Yang--Mills construction.
The original ERFL additionally queries the connected knowledge base during
preparation and persists the audit. These effects have no corresponding native
query or persistence occurrence in this two-stage projection. Original accepted,
quarantined, and stalled outcomes lead respectively to accept or stop, whereas the
native function returns $d$. Thus acceptance concerns an overlay disposition;
it must not silently acquire the meaning of a continuum existence proof.

% SPDX-FileCopyrightText: © 2026 Idiap Research Institute <contact@idiap.ch>
%
% SPDX-FileContributor: Andre Freitas <andre.freitas@idiap.ch>
% Neuro-symbolic AI Group
%
% SPDX-FileContributor: Danilo Gusicuma <danilo.gusicuma@idiap.ch>
% Neuro-symbolic AI Group
%
% SPDX-License-Identifier: MIT
% Generated from step-explanations.json
\paragraph{Step-by-step reconstruction.}
A selected observable shows exponential decay, while the proposed global spectral conclusion has no proof that the observable-generated sector is dense.
The following is a source-guided illustrative reconstruction. Internal procedures are prescriptions; the native declarations do not establish their execution.
\begin{enumerate}[leftmargin=*,label=\arabic*.]
\item \textbf{Prepare the continuum-boundary audit} (\texttt{continuum\_\allowbreak{}audit\_\allowbreak{}input}).
\begin{description}[style=standard,labelindent=0pt,leftmargin=5.5em,labelwidth=5em,labelsep=0.5em,align=left,font=\normalfont\bfseries,itemsep=0.15em,parsep=0pt,topsep=0.3em]
\item[Purpose] Bind the representation to its physical claim boundary.
\item[Input] Target context; representation proposal; continuum boundary
\item[Method] The native preparation is unresolved. Linked procedures specify carriers, observable maps, state transport, positivity, limit order and observable completeness as separate burdens.
\item[Check] Finite or sectorial evidence must remain distinguished from same-theory continuum and full-spectrum claims.
\item[Output] Intended audit input: target identity; boundary; candidate reference
\item[Next] Overlay assessment consumes it; query and persistence effects in original policy are not native occurrences.
\item[Example] Record decay for the visible observable family and mark density of its generated physical sector as an unproved assumption.
\end{description}
\item \textbf{Test the global inference} (\texttt{overlay\_\allowbreak{}disposition}).
\begin{description}[style=standard,labelindent=0pt,leftmargin=5.5em,labelwidth=5em,labelsep=0.5em,align=left,font=\normalfont\bfseries,itemsep=0.15em,parsep=0pt,topsep=0.3em]
\item[Purpose] Identify unsupported physical extrapolation.
\item[Input] Prepared candidate; retained continuum limits
\item[Method] Seek a sector invisible to the chosen observables and inspect whether the proposed transfer preserves the required structure.
\item[Check] Unchanged visible correlations cannot exclude low energies in an unobserved sector.
\item[Output] Overlay disposition: limited claim or quarantined global inference
\item[Next] Native output returns an overlay disposition; accepted status is not a continuum existence or mass-gap theorem.
\item[Example] In an abstract spectral calibration, adjoin an orthogonal sector with positive energies approaching zero. Visible decay is unchanged; this exposes the missing density inference without constructing a Yang--Mills counterexample.
\end{description}
\end{enumerate}
\paragraph{Changed-input contrast.}
If density and the required uniform extension estimates were independently established, the hidden-sector objection would need reassessment. Their mere appearance as named fields in an audit packet does not establish them.


\subsection{Deep review: challenge and repair the same epistemic state}
\label{app:granular-deep-review-and-reorientation}
The linked review skill makes disagreement itself an object of reasoning.
Reviewers inspect implemented code and persisted state through mathematical,
epistemic, conceptual, discovery and design lenses; they separate a mathematical
defect from an implementation defect and compare where the diagnoses converge or
conflict. Reflection then asks what the search changed in the understanding of
the problem: a new perspective, a sharper obstruction, persistent hardness or a
regression. A more elaborate artifact is not automatically a deeper result.
The named contract prescribes state formalization, advisory reviews, meta-review
and repair governance; this procedural ordering is a different layer from the
native parallel challenge/repair topology. Together they explain how a failure
becomes a revised research question while preserving unresolved disagreement,
rather than averaging reviewer opinions into apparent confirmation.

Deep review separates detecting why a research direction fails from proposing how
it might recover. Let $B,S,F,P,J,D$ abbreviate respectively
\texttt{Review\allowbreak Evidence\allowbreak Bundle}, \texttt{Formalized\allowbreak Review\allowbreak State},
\texttt{Review\allowbreak Defeater\allowbreak Assessment}, \texttt{Reorientation\allowbreak Repair\allowbreak Options},
\texttt{Reorientation\allowbreak Decision\allowbreak Input}, and \texttt{Reorientation\allowbreak Decision\allowbreak Assessment}.
\[
\begin{aligned}
&G:B\to\operatorname{Generated}[S],\quad
 R:S\to\operatorname{Generated}[F],\quad
 A:S\to\operatorname{Generated}[P],\\
&E:J\to\operatorname{Generated}[D],\qquad s=G(b),\\
&\left.\begin{aligned} f&=R(s.\mathrm{value})\\
                        p&=A(s.\mathrm{value})\end{aligned}\right\}
 \text{ declared parallel},\\
&j=\{\mathrm{defeater\_assessment}=f.\mathrm{value},\;
       \mathrm{repair\_options}=p.\mathrm{value}\},\qquad d=E(j).
\end{aligned}
\]
\texttt{review\_\allowbreak state} grounds the accumulated evidence in one formalized state.
\texttt{defeater\_\allowbreak assessment} attacks that state to identify bounded defects,
while \texttt{repair\_\allowbreak options} abductively proposes changes from the same state.
The two branches are complementary, not competing votes.
\texttt{Grounds\allowbreak Reorientation\allowbreak Alternatives} annotates both dependencies and makes
explicit that their point of departure is shared.
\texttt{reorientation\_\allowbreak decision} then calibrates the response using both named
fields of $J$. \texttt{Calibrates\allowbreak Reorientation\allowbreak Decision} records this joined
assessment boundary. The relation records express declared dependencies, not a
proof that a suggested repair resolves the defeater. Importantly, the repair
branch does not consume $f$: a repair proposed in parallel still has to be judged
against the defeater at the join.

For a constructed example, one review sees a promising correlation-decay
calculation, while another finds that the chosen observables miss a possible
low-energy sector. The disagreement is resolved by identifying its object: the
calculation may be valid for a visible sector while the global conclusion lacks
density. Reflection records a sharpening of the obstruction, not a completed gap
proof. Repair must construct a sufficiently rich observable class or restrict the
claim; simply collecting more decay curves would not address the defeater.
The output pairs a decision explanation with repair-assigned,
conceptual-reorientation-required, or stalled disposition. Original control routes
the first two to target navigation or conceptual hypothesis exploration and stops
on stall. Native reasoning returns the wrapped assessment; declared parallelism
and routing do not establish an observed execution.

% SPDX-FileCopyrightText: © 2026 Idiap Research Institute <contact@idiap.ch>
%
% SPDX-FileContributor: Andre Freitas <andre.freitas@idiap.ch>
% Neuro-symbolic AI Group
%
% SPDX-FileContributor: Danilo Gusicuma <danilo.gusicuma@idiap.ch>
% Neuro-symbolic AI Group
%
% SPDX-License-Identifier: MIT
% Generated from step-explanations.json
\paragraph{Step-by-step reconstruction.}
Reviewers examine a correlation-decay argument whose calculation concerns a visible sector but whose conclusion concerns the full physical spectrum.
The following is a source-guided illustrative reconstruction. Internal procedures are prescriptions; the native declarations do not establish their execution.
\begin{enumerate}[leftmargin=*,label=\arabic*.]
\item \textbf{Ground the reviewed state} (\texttt{review\_\allowbreak{}state}).
\begin{description}[style=standard,labelindent=0pt,leftmargin=5.5em,labelwidth=5em,labelsep=0.5em,align=left,font=\normalfont\bfseries,itemsep=0.15em,parsep=0pt,topsep=0.3em]
\item[Purpose] Ground criticism and repair in the same state.
\item[Input] Evidence bundle; implemented argument; persisted claims
\item[Method] Separate mathematical results, implementation behavior, assumptions and open obligations before interpreting progress.
\item[Check] Disagreement about scope must remain explicit rather than averaged into confidence.
\item[Output] Formalized review state: visible decay; unsupported global extension
\item[Next] Both parallel branches consume this same state; their coexistence gives no independence guarantee.
\item[Example] Write the state as a valid-looking visible-sector calculation plus the additional, unsupported claim that every physical excitation is gapped.
\end{description}
\item \textbf{Localize the defeater} (\texttt{defeater\_\allowbreak{}assessment}).
\begin{description}[style=standard,labelindent=0pt,leftmargin=5.5em,labelwidth=5em,labelsep=0.5em,align=left,font=\normalfont\bfseries,itemsep=0.15em,parsep=0pt,topsep=0.3em]
\item[Purpose] Locate the unsupported implication.
\item[Input] Grounded review state, shared with repair formation
\item[Method] Challenge quantifiers and representation coverage, separating a defect in the inference from a defect in implementation.
\item[Check] A useful diagnosis names the missing burden and the claim it affects.
\item[Output] Defeater assessment: missing observable-density premise
\item[Next] The join later compares this finding with repairs; the repair branch does not consume it directly.
\item[Example] An abstract direct-sum spectral countermodel preserves the reported correlations while adding energies approaching zero. It defeats the unrestricted correlation-to-gap inference and identifies the missing completeness premise; it is not a constructed Yang--Mills theory.
\end{description}
\item \textbf{Propose recovery options} (\texttt{repair\_\allowbreak{}options}).
\begin{description}[style=standard,labelindent=0pt,leftmargin=5.5em,labelwidth=5em,labelsep=0.5em,align=left,font=\normalfont\bfseries,itemsep=0.15em,parsep=0pt,topsep=0.3em]
\item[Purpose] Propose repairs that preserve useful content.
\item[Input] Formalized review state; no parallel refutation output
\item[Method] Propose a changed representation, missing bridge or narrower claim, explaining what each option would preserve.
\item[Check] A repair proposal must alter the defective dependency rather than relabel the current result.
\item[Output] Repair alternatives: richer observables or sectorial theorem
\item[Next] Final calibration must test these options against the separately produced defeater.
\item[Example] From the visible/global mismatch, propose either constructing a dense observable-generated subspace or explicitly restricting the conclusion to the visible sector.
\end{description}
\item \textbf{Judge repairs against criticism} (\texttt{reorientation\_\allowbreak{}decision}).
\begin{description}[style=standard,labelindent=0pt,leftmargin=5.5em,labelwidth=5em,labelsep=0.5em,align=left,font=\normalfont\bfseries,itemsep=0.15em,parsep=0pt,topsep=0.3em]
\item[Purpose] Match repairs to diagnosed defects.
\item[Input] Defeater assessment + parallel repair options
\item[Method] Check which proposal addresses the missing implication and record whether understanding sharpened, stalled or regressed.
\item[Check] A proposed repair does not resolve a defect merely by appearing in the same packet.
\item[Output] Reorientation assessment: repair assignment or conceptual reorientation
\item[Next] Original policy routes assignments; native return does not execute repair or prove density.
\item[Example] The narrower claim resolves the overstatement; the richer observable construction remains a research task. More curves without new sectors do not answer the defeater.
\end{description}
\end{enumerate}
\paragraph{Changed-input contrast.}
If a proposed repair only adds more decay curves from the same incomplete observable family, the join should reject it as irrelevant to density. A narrower sectorial claim changes the conclusion honestly.


\subsection{Contribution cultivation: test what a result enables}
\label{app:granular-scientific-contribution-cultivation}
The linked contribution skills separate finding a useful object from judging
its value and its effect on the target. Detection asks whether a subresult can
stand alone as a bridge, obstruction, reusable method or explanatory reframing,
then tests the smallest proposed claim against prior art, hidden premises and
proxy drift. Value assessment compares what would remain unknown without the
result, distinguishing exploration payoff from paper readiness. Alignment asks
whether the source object can feed the original target through a stated transfer;
enablement compares the precondition before the result with the postcondition
actually consumed by a next proof, falsifier or route decision. A beautiful local
calculation can therefore be valuable yet non-enabling, while a negative result
can enable a decisive pivot. The named contract explicitly separates these
assessments; they refine the meaning of the three native occurrences rather than
adding undocumented operations to that topology.

A supported result does not determine its scientific contribution by itself.
Contribution cultivation proposes an interpretation, challenges it, and assesses
readiness without inflating novelty. Let $S,H,T,J,Q$ abbreviate
\texttt{Result\allowbreak Atlas\allowbreak Slice}, \texttt{Scientific\allowbreak Contribution\allowbreak Hypothesis},
\texttt{Contribution\allowbreak Support\allowbreak Assessment}, \texttt{Contribution\allowbreak Readiness\allowbreak Input},
and \texttt{Contribution\allowbreak Readiness\allowbreak Assessment}:
\[
\begin{aligned}
&A:S\to\operatorname{Generated}[H],\qquad
 R:H\to\operatorname{Generated}[T],\qquad
 E:J\to\operatorname{Generated}[Q],\\
&h=A(s),\qquad t=R(h.\mathrm{value}),\\
&j=\{\mathrm{contribution\_hypothesis}=h.\mathrm{value},\;
       \mathrm{support\_assessment}=t.\mathrm{value}\},\qquad q=E(j).
\end{aligned}
\]
Occurrence \texttt{a0} abductively interprets the result as a candidate
contribution. Occurrence \texttt{r0} tests whether the result supports that
interpretation, under \texttt{Tests\allowbreak Contribution\allowbreak Hypothesis}. Occurrence
\texttt{e0} retains both interpretation and assessment at the named join and
calibrates readiness under \texttt{Calibrates\allowbreak Contribution\allowbreak Readiness}. The first
relation marks a test dependency; the second marks a calibration dependency.
Neither supplies missing evidence merely by being declared. The output contains
a disposition together with \texttt{Contribution\allowbreak Value\allowbreak Assessment},
\texttt{Main\allowbreak Problem\allowbreak Alignment\allowbreak Audit}, and \texttt{Enablement\allowbreak Audit}. These are returned
components, not three additional declared operation occurrences. Notice that
$q.\mathrm{value}$ unwraps \texttt{Generated}, whereas the assessment's own
\texttt{value} field denotes the contribution-value audit.

For a constructed example, cycle graphs show that connectedness alone cannot
support a spectral-gap lower bound uniform in graph size. An initial contribution
hypothesis might claim a general obstruction to fast diffusion. Testing rejects
that breadth: the result obstructs the uniform claim for the stated family and
assumptions. A defensible contribution can instead identify the missing structural
condition, align it with the uniform-decay problem, and enable a sharper search
for sufficient assumptions. Readiness still requires support for significance;
this illustration establishes neither novelty nor publication merit.
Original control stops when \texttt{r0} is refuted, routes paper-ready to
\texttt{manuscript\_\allowbreak result\_\allowbreak elaboration}, and stops on stall. The native sequence
always displays \texttt{e0} after \texttt{r0} and returns its wrapped assessment;
its handoff and stop comments are explicitly non-executable. A refuted
interpretation therefore must not be narrated as necessarily reaching the native
full-path calibration in the original policy.

% SPDX-FileCopyrightText: © 2026 Idiap Research Institute <contact@idiap.ch>
%
% SPDX-FileContributor: Andre Freitas <andre.freitas@idiap.ch>
% Neuro-symbolic AI Group
%
% SPDX-FileContributor: Danilo Gusicuma <danilo.gusicuma@idiap.ch>
% Neuro-symbolic AI Group
%
% SPDX-License-Identifier: MIT
% Generated from step-explanations.json
\paragraph{Step-by-step reconstruction.}
A completed cycle-graph example defeats a size-independent spectral gap from connectedness alone. The task is to identify what this negative result enables.
The following is a source-guided illustrative reconstruction. Internal procedures are prescriptions; the native declarations do not establish their execution.
\begin{enumerate}[leftmargin=*,label=\arabic*.]
\item \textbf{Interpret a possible contribution} (\texttt{a0}).
\begin{description}[style=standard,labelindent=0pt,leftmargin=5.5em,labelwidth=5em,labelsep=0.5em,align=left,font=\normalfont\bfseries,itemsep=0.15em,parsep=0pt,topsep=0.3em]
\item[Purpose] Identify the bounded result’s scientific use.
\item[Input] Result-atlas slice; assumptions; failure evidence
\item[Method] Propose the smallest standalone bridge, obstruction or reframing and identify who could use it.
\item[Check] A contribution hypothesis must not inherit the breadth of the original unsolved target.
\item[Output] Contribution hypothesis: obstruction to the insufficient premise
\item[Next] Refutation tests this interpretation, not merely the underlying calculation.
\item[Example] Propose that the cycle example shows why connectedness is too weak for the intended uniform diffusion bound and motivates a stronger structural condition.
\end{description}
\item \textbf{Challenge support and significance} (\texttt{r0}).
\begin{description}[style=standard,labelindent=0pt,leftmargin=5.5em,labelwidth=5em,labelsep=0.5em,align=left,font=\normalfont\bfseries,itemsep=0.15em,parsep=0pt,topsep=0.3em]
\item[Purpose] Test contribution scope and novelty support.
\item[Input] Proposed obstruction interpretation; stated scope
\item[Method] Test hidden premises, proxy drift and prior-art boundaries; compare what the result actually rules out.
\item[Check] Neither one slow family nor a true calculation establishes a general impossibility of fast diffusion.
\item[Output] Support assessment: bounded obstruction; unestablished breadth and novelty
\item[Next] Final readiness joins this finding with the same contribution hypothesis.
\item[Example] Reject the claim that fast diffusion is impossible in general; retain only failure of connectedness to imply a uniform gap over the specified family.
\end{description}
\item \textbf{Assess value and enablement} (\texttt{e0}).
\begin{description}[style=standard,labelindent=0pt,leftmargin=5.5em,labelwidth=5em,labelsep=0.5em,align=left,font=\normalfont\bfseries,itemsep=0.15em,parsep=0pt,topsep=0.3em]
\item[Purpose] Assess target value and communication readiness.
\item[Input] Contribution hypothesis + its support assessment
\item[Method] Separate conceptual value, target alignment and actual downstream enablement; compare the prior missing condition with the next task's requirements.
\item[Check] Exploration payoff cannot by itself establish paper readiness or novelty.
\item[Output] Readiness record: useful obstruction; open merit burden
\item[Next] Original policy can route to elaboration or stop; returned audit components are not extra native occurrences.
\item[Example] The result enables replacing an inadequate premise by a search for sufficient structural assumptions; any claim of publishable novelty still requires separate evidence.
\end{description}
\end{enumerate}
\paragraph{Changed-input contrast.}
If the result merely restates a known example without a new warranted use, its mathematical validity does not establish novelty or publication merit. Enablement needs a specific downstream change.


